% ============================================================================
%  Laboratorio di Programmazione -- Lezione 10 (10 novembre 2026)
%  Compilare con:  pdflatex -shell-escape Lezione10_Eccezioni_Ricorsione.tex
%  Richiede: mermaid.sty (nel progetto) + Mermaid CLI `mmdc` nel PATH
%  NB: ogni frame che contiene codice (lstlisting) o un diagramma (mermaid)
%      deve essere dichiarato [fragile].
% ============================================================================
\documentclass[10pt,aspectratio=169]{beamer}
\usetheme{Madrid}
\usecolortheme{default}

% `provide=*` carica la localizzazione italiana dai file .ini di babel: funziona anche
% su installazioni TeX prive di texlive-lang-italian; con una TeX Live completa e`
% equivalente al classico \usepackage[italian]{babel}.
\usepackage[italian,provide=*]{babel}
\usepackage[utf8]{inputenc}

% Localizzazione dei nomi dei blocchi Beamer
\uselanguage{italian}
\languagepath{italian}
\deftranslation[to=italian]{Definition}{Definizione}
\deftranslation[to=italian]{Theorem}{Teorema}
\deftranslation[to=italian]{Example}{Esempio}
\deftranslation[to=italian]{Proof}{Dimostrazione}
\deftranslation[to=italian]{Corollary}{Corollario}
\deftranslation[to=italian]{Lemma}{Lemma}

\usepackage{listings}
\usepackage{xcolor}
\usepackage{hyperref}
\usepackage{array}
\usepackage{booktabs}
\usepackage{mermaid}

% Mermaid: sfondo trasparente e dimensione massima adatta a una slide 16:9
\MermaidRendererOptions{-b transparent}
\MermaidAdjustBoxOpts{max width=0.95\linewidth,max height=0.62\textheight,center}

% Configurazione colori per il codice
\definecolor{codegreen}{rgb}{0,0.6,0}
\definecolor{codegray}{rgb}{0.5,0.5,0.5}
\definecolor{codepurple}{rgb}{0.58,0,0.82}
\definecolor{backcolour}{rgb}{0.95,0.95,0.92}
\definecolor{keywordblue}{rgb}{0.0,0.0,0.7}

% Stile per snippet Java
\lstdefinestyle{javastyle}{
    backgroundcolor=\color{backcolour},
    commentstyle=\color{codegreen}\itshape,
    keywordstyle=\color{keywordblue}\bfseries,
    numberstyle=\scriptsize\color{codegray},
    stringstyle=\color{codepurple},
    basicstyle=\ttfamily\footnotesize,
    breakatwhitespace=false,
    breaklines=true,
    captionpos=b,
    keepspaces=true,
    numbers=none,
    numbersep=5pt,
    showspaces=false,
    showstringspaces=false,
    showtabs=false,
    tabsize=2,
    aboveskip=0.6\medskipamount,
    belowskip=0.6\medskipamount,
    language=Java,
    morekeywords={default},
    literate={à}{{\`a}}1 {è}{{\`e}}1 {é}{{\'e}}1 {ì}{{\`i}}1
             {ò}{{\`o}}1 {ù}{{\`u}}1 {È}{{\`E}}1
}
\lstset{style=javastyle}

% Stile per righe di terminale
\lstdefinestyle{shell}{
    language=bash,
    backgroundcolor=\color{black!85},
    basicstyle=\ttfamily\footnotesize\color{white},
    keywordstyle=\color{white},
    stringstyle=\color{white},
    commentstyle=\color{green!70},
    numbers=none,
    frame=none,
    aboveskip=0.6\medskipamount,
    belowskip=0.6\medskipamount,
    showstringspaces=false
}

\title[Lezione 10 -- Eccezioni e ricorsione]{Lezione 10 -- Eccezioni e ricorsione}
\subtitle{\texttt{try}, \texttt{catch}, \texttt{finally}, \texttt{throw} e \texttt{throws}, eccezioni controllate e non, le nostre eccezioni; casi base e ricorsivi, lo stack delle chiamate}
\author[Mattia Fonisto]{Mattia Fonisto}
\institute[Lab.\ di Programmazione]{Laboratorio di Programmazione -- Ingegneria Informatica, II anno\\
           Università degli Studi di Napoli Federico II -- Sede di San Giovanni a Teduccio}
\date{10 novembre 2026}

\begin{document}

\frame{\titlepage}

% ========================================================================
\begin{frame}{Agenda della lezione (2h)}
\begin{enumerate}
    \item \textbf{Ripasso e Homework 9} \hfill (\emph{15 min})
    \begin{itemize}
        \item uno di voi alla lavagna con \texttt{Spedibile}, \texttt{Termostato} o \texttt{Magazzino}: lo commentiamo insieme
    \end{itemize}
    \item \textbf{Le eccezioni} \hfill (\emph{25 min})
    \begin{itemize}
        \item errori a tempo di esecuzione e stack trace; \texttt{try} e \texttt{catch}; \texttt{throw}
        \item propagazione e rilancio; la gerarchia, \texttt{throws}, eccezioni controllate e non; \texttt{finally}; le nostre eccezioni
    \end{itemize}
    \item \textbf{La ricorsione} \hfill (\emph{20 min})
    \begin{itemize}
        \item casi base e ricorsivi; il fattoriale; lo stack delle chiamate; terminazione
        \item ricorsione o iterazione; da ricorsivo a iterativo
    \end{itemize}
    \item \textbf{Esercitazione} -- tutti insieme \hfill (\emph{50 min})
    \begin{itemize}
        \item \emph{Commentiamo insieme}: tre metodi ricorsivi, letti riga per riga
        \item \emph{Insieme, passo passo}: il conto bancario che segnala gli errori con le eccezioni
    \end{itemize}
    \item \textbf{Wrap-up e Homework 10} \hfill (\emph{10 min})
\end{enumerate}
\end{frame}

% ========================================================================
\section{Le eccezioni}

\begin{frame}[fragile]{Eccezioni: errori a tempo di esecuzione -- \texttt{Divisione.java}}
\begin{columns}[T]
\column{0.56\textwidth}
\begin{lstlisting}[basicstyle=\ttfamily\scriptsize]
public class Divisione {

    public static int quoziente(int a, int b) {
        return a / b;
    }

    public static void main(String[] args) {
        Scanner in = new Scanner(System.in);
        System.out.print("Due interi: ");
        int a = in.nextInt();
        int b = in.nextInt();
        int q = quoziente(a, b);
        System.out.println("quoziente: " + q);
        System.out.println("fine");
    }
}
\end{lstlisting}
\begin{lstlisting}[style=shell,basicstyle=\ttfamily\scriptsize\color{white}]
Due interi: 7 0
Exception in thread "main"
java.lang.ArithmeticException: / by zero
    at Divisione.quoziente(Divisione.java:7)
    at Divisione.main(Divisione.java:15)
\end{lstlisting}

\column{0.44\textwidth}
\begin{block}{Eccezione}
\small Un'\emph{eccezione} è un oggetto che la JVM crea quando un'operazione non può essere completata. L'esecuzione si interrompe e l'eccezione è \textbf{sollevata} e rimandata a chi ha chiamato; se nessuno la gestisce, il programma termina con lo \emph{stack trace}.
\end{block}
\footnotesize
\begin{itemize}\itemsep=3pt
    \item già incontrate: \texttt{ClassCastException} (lezione~7), \texttt{NullPointerException} (lezione~8)
    \item \textbf{lo stack trace}: classe e messaggio (\texttt{/ by zero}), poi una riga per metodo, con file e numero di riga: in cima dove è nata (\texttt{quoziente}, riga 7 del file), sotto chi l'ha chiamata, fino al \texttt{main}
    \item \texttt{fine} non viene stampato: anche \texttt{main} è terminato
\end{itemize}
\end{columns}
\end{frame}

% ========================================================================
\begin{frame}[fragile]{Gestire un'eccezione: \texttt{try} e \texttt{catch} -- \texttt{Divisione.java} (versione 2)}
\begin{columns}[T]
\column{0.52\textwidth}
\begin{lstlisting}[basicstyle=\ttfamily\scriptsize]
try {
    int a = in.nextInt();
    int b = in.nextInt();
    int q = quoziente(a, b);
    System.out.println("quoziente: " + q);
} catch (InputMismatchException e) {
    System.out.println("servono due interi");
} catch (ArithmeticException e) {
    System.out.println(e.getMessage());
}
System.out.println("fine");
\end{lstlisting}
\begin{lstlisting}[style=shell,basicstyle=\ttfamily\scriptsize\color{white}]
Due interi: 7 0
/ by zero
fine
\end{lstlisting}
\begin{lstlisting}[style=shell,basicstyle=\ttfamily\scriptsize\color{white}]
Due interi: 7 x
servono due interi
fine
\end{lstlisting}
\footnotesize
\texttt{InputMismatchException} va importata da \texttt{java.util}, come \texttt{Scanner}; \texttt{ArithmeticException} sta in \texttt{java.lang}. \texttt{a}, \texttt{b} e \texttt{q} sono locali al blocco \texttt{try} (lezione~3): dopo non esistono più.

\column{0.48\textwidth}
\begin{block}{\texttt{try} e \texttt{catch}}
\small Nel blocco \texttt{try} sta il codice che può sollevare un'eccezione; ogni \texttt{catch} dichiara il \textbf{tipo} che sa gestire e un parametro, \texttt{e}, riferimento all'oggetto eccezione. Se nel \texttt{try} viene sollevata un'eccezione, le istruzioni che seguono nel \texttt{try} sono saltate e viene eseguito il \textbf{primo} \texttt{catch} del tipo giusto; poi si prosegue dopo l'ultimo \texttt{catch}. Se non viene sollevato nulla, i \texttt{catch} sono ignorati.
\end{block}
\footnotesize
\begin{itemize}\itemsep=3pt
    \item con \texttt{7 0}: la \texttt{println} del quoziente è saltata, si esegue il secondo \texttt{catch}, poi \texttt{fine}: il programma \textbf{non} termina. \texttt{e.getMessage()} è il messaggio; \texttt{println(e)} stampa classe e messaggio
\end{itemize}
\end{columns}
\end{frame}

% ========================================================================
\begin{frame}[fragile]{Sollevare un'eccezione: \texttt{throw} -- \texttt{Media.java}}
\begin{columns}[T]
\column{0.52\textwidth}
\begin{lstlisting}[basicstyle=\ttfamily\scriptsize]
/** La media di v, che non deve essere vuoto. */
public static double media(int[] v) {
    if (v.length == 0) {
        throw new IllegalArgumentException(
                "media di un array vuoto");
    }
    int somma = 0;
    for (int x : v) {
        somma += x;
    }
    return (double) somma / v.length;
}
\end{lstlisting}
\begin{block}{\texttt{throw}}
\small Un metodo che rileva una situazione che non sa gestire crea un oggetto eccezione con \texttt{new} e lo solleva con \texttt{throw}: il metodo \textbf{termina subito}, senza restituire nulla, e l'eccezione passa al chiamante. \textbf{Rileva} chi se ne può accorgere (il metodo); \textbf{gestisce} chi sa che cosa fare (il client).
\end{block}

\column{0.48\textwidth}
\begin{lstlisting}[basicstyle=\ttfamily\scriptsize]
public static void main(String[] args) {
    int[] voti = { 27, 30, 24 };
    int[] vuoto = {};
    System.out.println(media(voti));
    try {
        System.out.println(media(vuoto));
    } catch (IllegalArgumentException e) {
        System.out.println(e.getMessage());
    }
}
\end{lstlisting}
\begin{lstlisting}[style=shell,basicstyle=\ttfamily\scriptsize\color{white}]
27.0
media di un array vuoto
\end{lstlisting}
\footnotesize
\begin{itemize}\itemsep=2pt
    \item un \texttt{boolean} di esito non basta quando il metodo deve restituire un \textbf{valore}: nessun \texttt{double} può dire ``media di niente'' (\texttt{0.0} è una media legittima)
    \item \texttt{a / b} con \texttt{b} zero è come un \texttt{throw} di \texttt{ArithmeticException}, fatto dalla JVM. \texttt{IllegalArgumentException} sta in \texttt{java.lang}
\end{itemize}
\end{columns}
\end{frame}

% ========================================================================
\begin{frame}[fragile]{Propagazione: lo stack unwinding -- \texttt{Catena.java}}
\begin{columns}[T]
\column{0.55\textwidth}
\begin{lstlisting}[basicstyle=\ttfamily\scriptsize]
public static void main(String[] args) {
    primo();
}

public static void primo() {
    secondo();
}

public static void secondo() {
    terzo();
}

public static void terzo() {
    throw new IllegalArgumentException("errore");
}
\end{lstlisting}
\begin{lstlisting}[style=shell,basicstyle=\ttfamily\scriptsize\color{white}]
Exception in thread "main"
java.lang.IllegalArgumentException: errore
    at Catena.terzo(Catena.java:17)
    at Catena.secondo(Catena.java:13)
    at Catena.primo(Catena.java:9)
    at Catena.main(Catena.java:5)
\end{lstlisting}

\column{0.45\textwidth}
\MermaidAdjustBoxOpts{max width=\linewidth,max height=0.37\textheight,center}
\begin{mermaid}
---
config:
  themeVariables:
    fontSize: 32px
  flowchart:
    nodeSpacing: 14
    rankSpacing: 10
    wrappingWidth: 600
---
flowchart TB
  T["<b>terzo</b>: throw, termina"]
  S["<b>secondo</b>: nessun catch, termina"]
  P["<b>primo</b>: nessun catch, termina"]
  M["<b>main</b>: nessun catch, termina"]
  E["la JVM stampa lo stack trace"]
  T -- "l'eccezione risale" --> S --> P --> M --> E
  style T fill:#fee2e2,stroke:#b91c1c
  style E fill:#fef3c7,stroke:#b45309
\end{mermaid}
\MermaidAdjustBoxOpts{max width=0.95\linewidth,max height=0.62\textheight,center}
\begin{block}{Stack unwinding}
\small Un'eccezione non catturata \textbf{risale la catena delle chiamate}: ogni metodo attraversato termina e viene rimosso dallo stack, finché un \texttt{catch} la cattura; se non la cattura nessuno, nemmeno il \texttt{main}, il programma termina.
\end{block}
\end{columns}
\end{frame}

% ========================================================================
\begin{frame}[fragile]{Catturare lungo la catena e rilanciare -- \texttt{CatenaRilancio.java}}
\begin{columns}[T]
\column{0.54\textwidth}
\begin{lstlisting}[basicstyle=\ttfamily\scriptsize]
public static void main(String[] args) {
    primo();
    System.out.println("fine");
}

public static void primo() {
    try {
        secondo();
    } catch (IllegalArgumentException e) {
        System.out.println(e);
    }
}

public static void secondo() {
    try {
        terzo();
    } catch (IllegalArgumentException e) {
        System.out.println("secondo: rilancio");
        throw e;
    }
}
\end{lstlisting}

\column{0.46\textwidth}
\begin{lstlisting}[style=shell,basicstyle=\ttfamily\scriptsize\color{white}]
secondo: rilancio
java.lang.IllegalArgumentException: errore
fine
\end{lstlisting}
\footnotesize
\begin{itemize}\itemsep=3pt
    \item il \texttt{catch} può stare in \textbf{qualunque} metodo della catena: \texttt{primo} cattura ciò che nasce in \texttt{terzo}, due livelli più in basso
    \item la risalita si ferma al primo \texttt{catch} del tipo giusto; da lì l'esecuzione \textbf{prosegue normalmente}: \texttt{primo} stampa l'eccezione (classe e messaggio) e termina come sempre, \texttt{main} stampa \texttt{fine}
    \item \textbf{rilancio}: \texttt{secondo} cattura, fa la sua parte (un messaggio) e rilancia \textbf{lo stesso oggetto} con \texttt{throw e}: l'eccezione riprende a risalire. \texttt{terzo} è quello di prima: chi la solleva non sa, e non deve sapere, chi la gestirà
    \item il tipo nel \texttt{catch} può essere una \textbf{superclasse} di quella sollevata (upcasting, lezione~7): qui \texttt{catch (RuntimeException e)} funzionerebbe
\end{itemize}
\end{columns}
\end{frame}

% ========================================================================
\begin{frame}[fragile]{La gerarchia delle eccezioni: controllate e non controllate}
\begin{columns}[T]
\column{0.52\textwidth}
\MermaidAdjustBoxOpts{max width=\linewidth,max height=0.5\textheight,center}
\begin{mermaid}
---
config:
  themeVariables:
    fontSize: 28px
  class:
    hideEmptyMembersBox: true
---
classDiagram
  direction LR
  class Throwable
  class Error
  class Exception
  class RuntimeException
  class StackOverflowError
  class CloneNotSupportedException
  class ArithmeticException
  class IllegalArgumentException
  class NullPointerException
  class ClassCastException
  Throwable <|-- Error
  Throwable <|-- Exception
  Error <|-- StackOverflowError
  Exception <|-- CloneNotSupportedException
  Exception <|-- RuntimeException
  RuntimeException <|-- ArithmeticException
  RuntimeException <|-- IllegalArgumentException
  RuntimeException <|-- NullPointerException
  RuntimeException <|-- ClassCastException
\end{mermaid}
\MermaidAdjustBoxOpts{max width=0.95\linewidth,max height=0.62\textheight,center}
\footnotesize
Le eccezioni sono \textbf{oggetti di classi}, in una gerarchia (lezione~7) la cui radice è la classe \texttt{Throwable} (una classe, non un'interfaccia). Una parte: \texttt{InputMismatchException} e \texttt{ArrayIndexOutOfBoundsException} discendono da \texttt{RuntimeException}, \texttt{StackOverflowError} da \texttt{Error}, a più livelli.

\vspace{0.1cm}
Fra più \texttt{catch} vince il \textbf{primo} compatibile: la sottoclasse va scritta prima della superclasse.

\column{0.48\textwidth}
\begin{block}{Controllate e non controllate}
\small Le sottoclassi di \texttt{RuntimeException} sono \emph{non controllate} (\emph{unchecked}): di solito errori di programmazione, che si evitano scrivendo bene il codice. Le altre sottoclassi di \texttt{Exception} sono \emph{controllate} (\emph{checked}): situazioni da prevedere; il \textbf{compilatore controlla} che chi può riceverne una la gestisca o la dichiari.
\end{block}
\footnotesize
\begin{itemize}\itemsep=2pt
    \item \texttt{Error}: guasti gravi, come lo \emph{stack overflow} (con la ricorsione): non controllati, non si dovrebbero catturare
    \item un \texttt{catch} cattura anche le \textbf{sottoclassi} del suo tipo: \texttt{catch (Exception e)} prende tutte le eccezioni, anche quelle che non ci aspettavamo: meglio il tipo più specifico
\end{itemize}
\end{columns}
\end{frame}

% ========================================================================
\begin{frame}[fragile]{Dichiarare le eccezioni: \texttt{throws} -- \texttt{ProvaClone.java}}
\begin{columns}[T]
\column{0.5\textwidth}
{\footnotesize In \texttt{Persona implements Cloneable}, come nella lezione~8:}
\begin{lstlisting}[basicstyle=\ttfamily\scriptsize]
@Override
public Persona clone()
        throws CloneNotSupportedException {
    return (Persona) super.clone();
}
\end{lstlisting}
{\footnotesize Il client \textbf{cattura}, invece di dichiarare \texttt{throws} sul \texttt{main}:}
\begin{lstlisting}[basicstyle=\ttfamily\scriptsize]
public static void main(String[] args) {
    Persona p = new Persona("Anna", "Russo");
    try {
        Persona q = p.clone();
        System.out.println("copia: " + q);
    } catch (CloneNotSupportedException e) {
        System.out.println(e);
    }
}
\end{lstlisting}
\begin{lstlisting}[style=shell,basicstyle=\ttfamily\scriptsize\color{white}]
copia: Anna Russo
\end{lstlisting}
\footnotesize
Per le eccezioni non controllate \texttt{throws} non è obbligatorio: \texttt{media} non dichiara la sua.

\column{0.5\textwidth}
\begin{block}{\texttt{throws}}
\small Nell'intestazione, \texttt{throws} elenca le eccezioni \textbf{controllate} che possono uscire dal metodo, sollevate da lui stesso o da un metodo che chiama, e che il metodo non cattura: le \textbf{rimanda al chiamante}. Chi lo chiama, a sua volta, deve \textbf{catturarle} (\texttt{try}/\texttt{catch}) oppure \textbf{rimandarle al proprio chiamante}, dichiarandole con \texttt{throws}.
\end{block}
\footnotesize
\begin{itemize}\itemsep=2pt
    \item finalmente la riga della lezione~8: \texttt{Object.clone()} dichiara \texttt{throws CloneNotSupportedException} e la solleva se l'oggetto non è \texttt{Cloneable} (lezione~9). Senza \texttt{throws} né \texttt{catch}: \emph{``unreported exception CloneNotSupportedException; must be caught or declared''}
    \item senza \texttt{implements Cloneable} si compila lo stesso, ma si arriva al \texttt{catch}:
\end{itemize}
\begin{lstlisting}[style=shell,basicstyle=\ttfamily\scriptsize\color{white}]
java.lang.CloneNotSupportedException: Persona
\end{lstlisting}
\end{columns}
\end{frame}

% ========================================================================
\begin{frame}[fragile]{\texttt{finally} e multi-catch -- \texttt{DivisioneFinally.java}}
\begin{columns}[T]
\column{0.52\textwidth}
\begin{lstlisting}[basicstyle=\ttfamily\scriptsize]
try {
    int a = in.nextInt();
    int b = in.nextInt();
    System.out.println(a / b);
} catch (InputMismatchException
        | ArithmeticException e) {
    System.out.println("errore: " + e);
} finally {
    System.out.println("tentativo concluso");
}
System.out.println("fine");
\end{lstlisting}
\begin{lstlisting}[style=shell,basicstyle=\ttfamily\scriptsize\color{white}]
Due interi: 7 0
errore: java.lang.ArithmeticException: / by zero
tentativo concluso
fine
\end{lstlisting}
\begin{lstlisting}[style=shell,basicstyle=\ttfamily\scriptsize\color{white}]
Due interi: 7 2
3
tentativo concluso
fine
\end{lstlisting}

\column{0.48\textwidth}
\begin{block}{\texttt{finally}}
\small Il blocco \texttt{finally}, dopo i \texttt{catch}, viene eseguito \textbf{sempre}: che il \texttt{try} finisca normalmente, che un \texttt{catch} gestisca l'eccezione, che l'eccezione sia rimandata al chiamante o che il \texttt{try} esegua un \texttt{return}.
\end{block}
\footnotesize
\begin{itemize}\itemsep=3pt
    \item serve per il lavoro da fare \textbf{comunque}: ad esempio chiudere un file (lezione~13)
    \item \textbf{multi-catch}: un solo \texttt{catch} per più tipi separati da \texttt{|}, ammesso che abbiano una gestione comune. \texttt{e} ha il tipo dell'antenato comune più vicino: qui \texttt{RuntimeException}
    \item l'unico vincolo è fra i tipi \textbf{elencati}: nessuno può essere sottoclasse di un altro dell'elenco (\texttt{ArithmeticException | RuntimeException} non compila). Avere antenati in comune non conta: \texttt{StackOverflowError | InputMismatchException} è lecito, ed \texttt{e} è un \texttt{Throwable}
\end{itemize}
\end{columns}
\end{frame}

% ========================================================================
\begin{frame}[fragile]{Le nostre eccezioni -- \texttt{ParcheggioPienoException.java}}
\begin{columns}[T]
\column{0.5\textwidth}
\begin{lstlisting}[basicstyle=\ttfamily\scriptsize]
/** Sollevata da entra(): parcheggio pieno. */
public class ParcheggioPienoException
        extends Exception {

    public ParcheggioPienoException() {
        super("parcheggio pieno");
    }
}
\end{lstlisting}
{\footnotesize In \texttt{Parcheggio} (\texttt{nome}, \texttt{posti}, \texttt{occupati}):}
\begin{lstlisting}[basicstyle=\ttfamily\scriptsize]
/** Occupa un posto libero. */
public void entra()
        throws ParcheggioPienoException {
    if (occupati == posti) {
        throw new ParcheggioPienoException();
    }
    occupati++;
}
\end{lstlisting}
\footnotesize
Controllata: \texttt{p.entra()} senza \texttt{try} né \texttt{throws} non compila: \emph{``unreported exception ParcheggioPienoException; must be caught or declared to be thrown''}.

\column{0.5\textwidth}
\begin{lstlisting}[basicstyle=\ttfamily\scriptsize]
Parcheggio p = new Parcheggio("P1", 2);
try {
    p.entra();
    p.entra();
    p.entra();
    System.out.println("tre auto dentro");
} catch (ParcheggioPienoException e) {
    System.out.println(e.getMessage());
}
System.out.println(p);
\end{lstlisting}
\begin{lstlisting}[style=shell,basicstyle=\ttfamily\scriptsize\color{white}]
parcheggio pieno
P1: 2/2
\end{lstlisting}
\footnotesize
\begin{itemize}\itemsep=2pt
    \item un'eccezione nostra è una \textbf{sottoclasse} di \texttt{Exception} (controllata) o di \texttt{RuntimeException} (non controllata); il nome finisce per convenzione in \texttt{Exception}
    \item il costruttore passa il messaggio con \texttt{super(...)} (lezione~7): è quello che restituisce \texttt{getMessage()}. Può avere variabili di istanza e metodi, come ogni classe: nell'esercitazione
\end{itemize}
\end{columns}
\end{frame}

% ========================================================================
\begin{frame}[fragile]{Quando (non) usare le eccezioni -- \texttt{Quoziente.java}}
\begin{columns}[T]
\column{0.5\textwidth}
\begin{lstlisting}[basicstyle=\ttfamily\scriptsize]
int a = 7;
int b = 0;
// NO: l'eccezione usata come un if
try {
    System.out.println(a / b);
} catch (ArithmeticException e) {
    System.out.println("divisore nullo");
}
// SI': il controllo normale, con un if
if (b != 0) {
    System.out.println(a / b);
} else {
    System.out.println("divisore nullo");
}
\end{lstlisting}
\footnotesize
Le due versioni stampano lo stesso messaggio; la seconda dice nell'\texttt{if} qual è il caso particolare, la prima lo nasconde in un \texttt{catch}.

\column{0.5\textwidth}
\footnotesize
\begin{itemize}\itemsep=4pt
    \item le eccezioni servono quando \textbf{chi rileva} il problema non sa come gestirlo e \textbf{chi sa} gestirlo non sa come rilevarlo: \texttt{media} sa che l'array è vuoto, ma solo il client sa se avvisare, ripetere o rinunciare
    \item \textbf{non} al posto del normale controllo di flusso: se basta un \texttt{if} (e un \texttt{boolean}), si usa l'\texttt{if}. Prima di dividere si controlla il divisore; prima di un cast, \texttt{instanceof} (lezione~8)
    \item \textbf{controllata o non controllata?} Controllata se il chiamante \textbf{deve} prevedere il caso e può rimediare (il parcheggio pieno: cercarne un altro). Non controllata se il chiamante l'avrebbe evitata scrivendo bene il codice: una precondizione violata, come l'array vuoto passato a \texttt{media}
    \item un'eccezione non gestita termina il programma: meglio pensarci prima
\end{itemize}
\end{columns}
\end{frame}

% ========================================================================
\section{La ricorsione}

\begin{frame}[fragile]{Ricorsione: definire qualcosa in termini di se stesso}
\begin{columns}[T]
\column{0.52\textwidth}
\begin{block}{Ricorsione}
\small Una definizione è \emph{ricorsiva} quando descrive una cosa usando la cosa stessa, applicata a un caso \textbf{più piccolo}: la definizione si riapplica al caso più piccolo, e poi a quello ancora più piccolo, finché si arriva a un caso così semplice da avere la risposta \textbf{diretta}. Un metodo è ricorsivo se, per risolvere un problema, chiama se stesso su un problema dello stesso tipo ma più piccolo.
\end{block}
\footnotesize
\begin{itemize}\itemsep=3pt
    \item il \textbf{fattoriale}: $0! = 1$; per $n \geq 1$, $n! = n \cdot (n-1)!$: il fattoriale di $n$ è definito con il fattoriale di $n-1$
    \item una \textbf{stringa} è la stringa vuota, oppure un carattere seguito da una stringa: \texttt{"java"} è \texttt{'j'} seguito da \texttt{"ava"} (\texttt{s.charAt(0)} e \texttt{s.substring(1)}, lezione~2)
    \item cercare una parola nel \textbf{dizionario}: si apre a metà e si cerca nella metà giusta, che è ancora un dizionario, più piccolo (la ricerca binaria, lezione~11)
\end{itemize}

\column{0.48\textwidth}
\footnotesize
\textbf{Come si svolge un calcolo ricorsivo: $3!$}
\begin{enumerate}\itemsep=2pt
    \item \textbf{si scende}, applicando la definizione finché compare un caso con risposta diretta:\\
    $3! = 3 \cdot 2!$ \quad $2! = 2 \cdot 1!$ \quad $1! = 1 \cdot 0!$ \quad $0! = 1$
    \item \textbf{si risale}, sostituendo i valori ottenuti, dall'ultimo al primo:\\
    $1! = 1 \cdot 1 = 1$ \quad $2! = 2 \cdot 1 = 2$ \quad $3! = 3 \cdot 2 = 6$
\end{enumerate}
Ogni passo in discesa lascia in sospeso un calcolo ($3 \cdot {?}$), che si completa in risalita. Da qui gli ingredienti di un metodo ricorsivo:
\begin{itemize}\itemsep=2pt
    \item uno o più \textbf{casi base}: la risposta è diretta, senza chiamate ($0! = 1$; la stringa vuota)
    \item il \textbf{caso ricorsivo}: una chiamata su un problema \textbf{più piccolo}, il cui risultato si combina con il resto ($n \cdot (n-1)!$)
    \item la \textbf{convergenza}: ogni chiamata si avvicina a un caso base, altrimenti la discesa non finisce
\end{itemize}
\end{columns}
\end{frame}

% ========================================================================
\begin{frame}[fragile]{La ricorsione in Java: il fattoriale -- \texttt{Fattoriale.java}}
\begin{columns}[T]
\column{0.5\textwidth}
\begin{lstlisting}[basicstyle=\ttfamily\scriptsize]
/** n!, con 0 <= n <= 12. */
public static int fattoriale(int n) {
    if (n < 0) {
        throw new IllegalArgumentException(
                "n negativo: " + n);
    }
    if (n == 0) {
        return 1;       // caso base
    }
    // caso ricorsivo
    return n * fattoriale(n - 1);
}

public static void main(String[] args) {
    System.out.println(fattoriale(3));
    System.out.println(fattoriale(12));
}
\end{lstlisting}
\begin{lstlisting}[style=shell,basicstyle=\ttfamily\scriptsize\color{white}]
6
479001600
\end{lstlisting}
\footnotesize
$13! = 6\,227\,020\,800$ non sta in un \texttt{int}: da qui la precondizione $n \leq 12$. Un solo codice, tante attivazioni: come un metodo che ne chiama un altro.

\column{0.5\textwidth}
\MermaidAdjustBoxOpts{max width=\linewidth,max height=0.46\textheight,center}
\begin{mermaid}
---
config:
  themeVariables:
    fontSize: 30px
  flowchart:
    nodeSpacing: 20
    rankSpacing: 16
    wrappingWidth: 900
---
flowchart TB
  M["<b>main</b>: fattoriale(3)"]
  F3["<b>n = 3</b>: 3 * fattoriale(2) = 6"]
  F2["<b>n = 2</b>: 2 * fattoriale(1) = 2"]
  F1["<b>n = 1</b>: 1 * fattoriale(0) = 1"]
  F0["<b>n = 0</b>: caso base = 1"]
  M --> F3 --> F2 --> F1 --> F0
  style F0 fill:#dcfce7,stroke:#15803d
\end{mermaid}
\MermaidAdjustBoxOpts{max width=0.95\linewidth,max height=0.62\textheight,center}
\footnotesize
\begin{itemize}\itemsep=2pt
    \item \textbf{fase discendente} (dall'alto): ogni chiamata aspetta la successiva; lo stack cresce a ogni chiamata, ciascuna con il \textbf{suo} \texttt{n} (lezione~3)
    \item \textbf{fase ascendente} (dal basso): il caso base restituisce 1; ogni chiamata moltiplica, restituisce e viene rimossa dallo stack
\end{itemize}
\end{columns}
\end{frame}

% ========================================================================
\begin{frame}[fragile]{Terminazione e \texttt{StackOverflowError} -- \texttt{SenzaBase.java}}
\begin{columns}[T]
\column{0.5\textwidth}
\begin{lstlisting}[basicstyle=\ttfamily\scriptsize]
/** Nessun controllo su n < 0. */
public static int fattoriale(int n) {
    if (n == 0) {
        return 1;
    }
    return n * fattoriale(n - 1);
}

public static void main(String[] args) {
    System.out.println(fattoriale(-1));
}
\end{lstlisting}
\begin{lstlisting}[style=shell,basicstyle=\ttfamily\scriptsize\color{white}]
Exception in thread "main"
java.lang.StackOverflowError
    at SenzaBase.fattoriale(SenzaBase.java:9)
    at SenzaBase.fattoriale(SenzaBase.java:9)
    at SenzaBase.fattoriale(SenzaBase.java:9)
    ... (1024 righe at in tutto)
\end{lstlisting}

\column{0.5\textwidth}
\begin{block}{Terminazione}
\small Un metodo ricorsivo termina se ogni catena di chiamate arriva a un caso base: a ogni chiamata il problema deve \textbf{diminuire}, verso un caso base che prima o poi viene raggiunto. Il numero di chiamate annidate è la \emph{profondità} della ricorsione.
\end{block}
\footnotesize
\begin{itemize}\itemsep=3pt
    \item con $n = -1$: $-2$, $-3$, \ldots: $n$ diminuisce ma \textbf{si allontana} da 0. Ogni chiamata occupa spazio sullo stack, che ha una dimensione \textbf{finita}: quando è pieno, la JVM solleva \texttt{StackOverflowError}, un \texttt{Error}: non va catturato, va corretto il codice
    \item è l'equivalente ricorsivo del ciclo infinito; con un ciclo, però, il programma resta bloccato, qui termina con un errore
    \item il rimedio è nella versione della slide precedente: la precondizione $n \geq 0$ controllata con un \texttt{throw}
\end{itemize}
\end{columns}
\end{frame}

% ========================================================================
\begin{frame}[fragile]{Ricorsione o iterazione? -- \texttt{ContoAllaRovescia.java}}
\begin{columns}[T]
\column{0.5\textwidth}
\begin{lstlisting}[basicstyle=\ttfamily\scriptsize]
public static void contoAllaRovescia(int n) {
    if (n == 0) {
        System.out.println("via!");
    } else {
        System.out.println(n);
        contoAllaRovescia(n - 1);
    }
}

public static void contoIterativo(int n) {
    while (n != 0) {
        System.out.println(n);
        n = n - 1;
    }
    System.out.println("via!");
}
\end{lstlisting}
\footnotesize
Con \texttt{3}, entrambi stampano \texttt{3}, \texttt{2}, \texttt{1}, \texttt{via!}

\column{0.5\textwidth}
\footnotesize
\begin{itemize}\itemsep=4pt
    \item lo stesso problema, risolto in due modi: nella versione ricorsiva il ``ripeti'' è la chiamata a se stesso e l'arresto è il caso base; in quella iterativa il ``ripeti'' è il \texttt{while} e l'arresto è la sua condizione
    \item ogni metodo ricorsivo ha una versione iterativa: qui basta un \texttt{while}, per il fattoriale un \texttt{for} che accumula il prodotto; per altri problemi (la Torre di Hanoi dell'homework) è molto più complicata
    \item \textbf{quando usarla}: se il problema è ricorsivo per natura e la soluzione ricorsiva è più semplice e leggibile; l'iterazione se il ciclo è ovvio e le prestazioni contano
\end{itemize}
\vspace{0.1cm}
\begin{center}
\renewcommand{\arraystretch}{1.05}
\scriptsize
\begin{tabular}{@{}lll@{}}
\toprule
 & \textbf{iterazione} & \textbf{ricorsione} \\
\midrule
costrutto & \texttt{while}, \texttt{for} & \texttt{if} + chiamata \\
si ferma & condizione del ciclo & caso base \\
se non si ferma & ciclo infinito & \texttt{StackOverflowError} \\
costo & nessuna chiamata & una chiamata per passo \\
\bottomrule
\end{tabular}
\end{center}
\end{columns}
\end{frame}

% ========================================================================
\section{Esercitazione}

\begin{frame}[fragile]{Commentiamo insieme (20 min) -- il problema}
\begin{columns}[T]
\column{0.55\textwidth}
\footnotesize
\textbf{Il problema.} Tre piccoli problemi da risolvere \textbf{senza cicli}, solo con la ricorsione. Il codice è già scritto, nella classe \texttt{Ricorsivi}: lo leggiamo riga per riga cercando, in ogni metodo, il caso base, il caso ricorsivo e perché la ricorsione termina.

\vspace{0.1cm}
\textbf{I tre problemi}
\begin{enumerate}\itemsep=3pt
    \item la \textbf{somma delle cifre} di un intero non negativo: per \texttt{1234}, $1 + 2 + 3 + 4 = 10$. L'ultima cifra è \texttt{n \% 10}, il numero senza l'ultima cifra è \texttt{n / 10}
    \item la \textbf{stampa inversa} di un array: per \texttt{\{ 3, 8, 1, 6 \}}, \texttt{6 1 8 3}, senza indici che scendono
    \item i \textbf{numeri di Fibonacci}: $F_0 = 0$, $F_1 = 1$ e ciascuno dei successivi è la somma dei due precedenti: 0, 1, 1, 2, 3, 5, 8, 13, \ldots{} Il programma conta anche quante chiamate servono per calcolare $F_{20}$
\end{enumerate}

\column{0.45\textwidth}
\footnotesize
\textbf{Le chiamate per calcolare $F_4$}
\MermaidAdjustBoxOpts{max width=\linewidth,max height=0.6\textheight,center}
\begin{mermaid}
---
config:
  themeVariables:
    fontSize: 20px
  flowchart:
    nodeSpacing: 14
    rankSpacing: 30
    wrappingWidth: 600
---
flowchart TB
  A["fib(4)"] --> B["fib(3)"]
  A --> C["fib(2)"]
  B --> D["fib(2)"]
  B --> E["fib(1)"]
  D --> F["fib(1)"]
  D --> G["fib(0)"]
  C --> H["fib(1)"]
  C --> I["fib(0)"]
  style C fill:#fef3c7,stroke:#b45309
  style D fill:#fef3c7,stroke:#b45309
  style E fill:#dcfce7,stroke:#15803d
  style F fill:#dcfce7,stroke:#15803d
  style G fill:#dcfce7,stroke:#15803d
  style H fill:#dcfce7,stroke:#15803d
  style I fill:#dcfce7,stroke:#15803d
\end{mermaid}
\MermaidAdjustBoxOpts{max width=0.95\linewidth,max height=0.62\textheight,center}
In verde i casi base; in giallo lo stesso calcolo, fatto due volte.
\end{columns}
\end{frame}

% ========================================================================
\begin{frame}[fragile]{Commentiamo insieme -- \texttt{sommaCifre} e \texttt{stampaInversa}}
\begin{columns}[T]
\column{0.46\textwidth}
{\footnotesize Nella classe \texttt{Ricorsivi}:}
\begin{lstlisting}[basicstyle=\ttfamily\scriptsize]
private static int chiamate = 0;

/** La somma delle cifre di n, n >= 0. */
public static int sommaCifre(int n) {
    if (n < 10) {
        return n;
    }
    return n % 10 + sommaCifre(n / 10);
}
\end{lstlisting}
\footnotesize
\begin{itemize}\itemsep=2pt
    \item caso base: un numero di una cifra; caso ricorsivo: l'ultima cifra più la somma delle cifre del resto
    \item \texttt{sommaCifre(1234)} $=$ \texttt{4 + sommaCifre(123)} $=$ \texttt{4 + 3 + sommaCifre(12)} $=$ \texttt{4 + 3 + 2 + 1}
    \item \texttt{n / 10} ha una cifra in meno: dopo al più dieci chiamate si arriva sotto 10
\end{itemize}

\column{0.54\textwidth}
\begin{lstlisting}[basicstyle=\ttfamily\scriptsize]
/** Stampa v[i], v[i+1], ... in ordine inverso. */
public static void stampaInversa(int[] v, int i) {
    if (i < v.length) {
        stampaInversa(v, i + 1);
        System.out.print(v[i] + " ");
    }
}

/** Stampa tutto v al contrario: parte da 0. */
public static void stampaInversa(int[] v) {
    stampaInversa(v, 0);
    System.out.println();
}
\end{lstlisting}
\footnotesize
\begin{itemize}\itemsep=2pt
    \item il problema ``più piccolo'' è la parte di array da \texttt{i} in poi: \texttt{i} cresce fino a \texttt{v.length}, dove la parte è vuota e non c'è nulla da fare (caso base \textbf{implicito}: l'\texttt{if} è falso)
    \item prima la chiamata, \textbf{poi} la stampa: \texttt{v[0]} viene stampato per ultimo, quando tutte le chiamate successive sono già tornate
    \item due metodi in \textbf{overloading} (lezione~3): il client chiama quello senza indice
\end{itemize}
\end{columns}
\end{frame}

% ========================================================================
\begin{frame}[fragile]{Commentiamo insieme -- \texttt{fibonacci} e \texttt{main}}
\begin{columns}[T]
\column{0.52\textwidth}
\begin{lstlisting}[basicstyle=\ttfamily\scriptsize]
/** L'n-esimo numero di Fibonacci. */
public static int fibonacci(int n) {
    chiamate++;
    if (n <= 1) {
        return n;
    }
    return fibonacci(n - 1) + fibonacci(n - 2);
}

public static void main(String[] args) {
    System.out.println(sommaCifre(1234));
    int[] v = { 3, 8, 1, 6 };
    stampaInversa(v);
    System.out.println(fibonacci(20));
    System.out.println(chiamate + " chiamate");
}
\end{lstlisting}
\begin{lstlisting}[style=shell,basicstyle=\ttfamily\scriptsize\color{white}]
10
6 1 8 3
6765
21891 chiamate
\end{lstlisting}

\column{0.48\textwidth}
\footnotesize
\textbf{Rispondiamo insieme:}
\begin{enumerate}\itemsep=3pt
    \item in \texttt{sommaCifre}, il caso base potrebbe essere \texttt{n == 0}, con \texttt{return 0}? Che cosa cambierebbe?
    \item in \texttt{stampaInversa}, che cosa stamperebbe il metodo scambiando la chiamata ricorsiva e la \texttt{print}?
    \item \texttt{fibonacci} ha \textbf{due} casi base in un solo \texttt{if}: quali? Perché ne servono due?
    \item per $F_{20}$ servono 21\,891 chiamate: perché così tante? Guardate l'albero di $F_4$. Come lo calcolereste con un ciclo e due variabili?
    \item perché \texttt{chiamate} è una variabile \texttt{static} della classe e non una variabile locale di \texttt{fibonacci}?
\end{enumerate}
\end{columns}
\end{frame}

% ========================================================================
\begin{frame}[fragile]{Insieme, passo passo (30 min) -- il problema}
\begin{columns}[T]
\column{0.6\textwidth}
\footnotesize
\textbf{Il problema.} Un conto bancario deve \textbf{rifiutare} le operazioni impossibili e dire al client \textbf{perché}: un importo non positivo è un errore di chi scrive il programma; un prelievo più grande del saldo è una situazione normale, che il programma deve prevedere, e il client deve poter sapere \textbf{quanto manca}. Scriviamo \texttt{ContoBancario} e l'eccezione \texttt{SaldoInsufficienteException}. Il programma di prova \texttt{ProvaConto} è già scritto: legge l'importo da prelevare dalla tastiera, insistendo finché non arriva un numero positivo.

\vspace{0.1cm}
\textbf{Le regole}
\begin{itemize}\itemsep=2pt
    \item un conto ha un intestatario (che non cambia) e un saldo, che parte da 0; si leggono con due \texttt{get}
    \item \texttt{versa(importo)} e \texttt{preleva(importo)}: con un importo non positivo, nessuna modifica e un'eccezione con il messaggio \texttt{non positivo: -5.0}
    \item \texttt{preleva} con un importo maggiore del saldo: nessuna modifica e un'eccezione che dice quanto manca
    \item \texttt{trasferisci(altro, importo)}: un prelievo da questo conto e un versamento su \texttt{altro}; se il prelievo fallisce, nessun versamento
    \item \texttt{toString}: \texttt{Anna: saldo 100.0}
\end{itemize}

\column{0.4\textwidth}
\MermaidAdjustBoxOpts{max width=\linewidth,max height=0.66\textheight,center}
\begin{mermaid}
---
config:
  themeVariables:
    fontSize: 20px
  flowchart:
    nodeSpacing: 14
    rankSpacing: 30
    wrappingWidth: 600
---
flowchart TB
  P["<b>ProvaConto.java</b><br/>già scritto"]
  C["<b>ContoBancario.java</b><br/>da scrivere"]
  E["<b>SaldoInsufficienteException.java</b><br/>da scrivere"]
  P -- "usa" --> C
  C -- "solleva" --> E
  P -- "cattura" --> E
  style P fill:#dcfce7,stroke:#15803d
  style C fill:#fef3c7,stroke:#b45309
  style E fill:#fef3c7,stroke:#b45309
\end{mermaid}
\MermaidAdjustBoxOpts{max width=0.95\linewidth,max height=0.62\textheight,center}
\end{columns}
\end{frame}

% ========================================================================
\begin{frame}[fragile]{Insieme, passo passo -- \texttt{ProvaConto.java} (1/2): l'input robusto}
\begin{columns}[T]
\column{0.56\textwidth}
\begin{lstlisting}[basicstyle=\ttfamily\scriptsize]
/** Legge un importo positivo: insiste. */
public static double leggiImporto(Scanner in) {
    double importo = 0;
    while (importo <= 0) {
        System.out.print("Importo: ");
        try {
            importo = in.nextDouble();
        } catch (InputMismatchException e) {
            in.nextLine();   // scarta l'input
            System.out.println("non e' un numero");
        }
    }
    return importo;
}
\end{lstlisting}
\begin{lstlisting}[style=shell,basicstyle=\ttfamily\scriptsize\color{white}]
Importo: abc
non e' un numero
Importo: -5
Importo: 250
\end{lstlisting}

\column{0.44\textwidth}
\footnotesize
\begin{itemize}\itemsep=4pt
    \item \textbf{input robusto}: un programma che legge dalla tastiera non deve terminare con uno stack trace perché l'utente ha scritto una parola al posto di un numero; deve chiedere di nuovo
    \item se \texttt{nextDouble} trova \texttt{abc}, solleva \texttt{InputMismatchException} e \textbf{lascia} \texttt{abc} nell'input: \texttt{nextLine()} lo scarta. Senza quella riga il ciclo rileggerebbe \texttt{abc} all'infinito
    \item \texttt{importo} resta a 0 e il ciclo ripete; un numero non positivo, come \texttt{-5}, è letto ma non fa uscire dal ciclo: qui basta la condizione del \texttt{while}, senza eccezioni
    \item il \texttt{catch} sta \textbf{dentro} il ciclo: dopo un errore si riprova; fuori dal ciclo, al primo errore il ciclo finirebbe
\end{itemize}
\end{columns}
\end{frame}

% ========================================================================
\begin{frame}[fragile]{Insieme, passo passo -- \texttt{ProvaConto.java} (2/2): il \texttt{main}}
\begin{columns}[T]
\column{0.57\textwidth}
\begin{lstlisting}[basicstyle=\ttfamily\scriptsize]
Scanner in = new Scanner(System.in);
ContoBancario anna = new ContoBancario("Anna");
ContoBancario luca = new ContoBancario("Luca");
anna.versa(100);
try {
    anna.preleva(leggiImporto(in));
    System.out.println("prelievo eseguito");
} catch (SaldoInsufficienteException e) {
    System.out.println("mancano " + e.getMancante());
} finally {
    System.out.println(anna);
}
try {
    anna.trasferisci(luca, 70);
} catch (SaldoInsufficienteException e) {
    System.out.println(e.getMessage());
}
System.out.println(anna + ", " + luca);
\end{lstlisting}

\column{0.43\textwidth}
{\footnotesize Con \texttt{abc}, \texttt{-5}, \texttt{250} (dopo i prompt):}
\begin{lstlisting}[style=shell,basicstyle=\ttfamily\scriptsize\color{white}]
mancano 150.0
Anna: saldo 100.0
Anna: saldo 30.0, Luca: saldo 70.0
\end{lstlisting}
{\footnotesize Con \texttt{40}:}
\begin{lstlisting}[style=shell,basicstyle=\ttfamily\scriptsize\color{white}]
Importo: 40
prelievo eseguito
Anna: saldo 60.0
saldo insufficiente, mancano 10.0 euro
Anna: saldo 60.0, Luca: saldo 0.0
\end{lstlisting}
\footnotesize
\begin{itemize}\itemsep=2pt
    \item il primo \texttt{catch} usa \texttt{getMancante()}, un metodo della \textbf{nostra} eccezione; il \texttt{finally} stampa il conto in entrambi i casi
    \item \texttt{IllegalArgumentException} non viene catturata: un importo non positivo è un errore del programma, che si evita a monte (qui \texttt{leggiImporto} accetta solo importi positivi), non si gestisce
\end{itemize}
\end{columns}
\end{frame}

% ========================================================================
\begin{frame}[fragile]{Insieme, passo passo -- le decisioni}
\begin{columns}[T]
\column{0.62\textwidth}
\footnotesize
\textbf{Prima di scrivere il codice, decidiamo insieme:}
\begin{enumerate}\itemsep=3pt
    \item \textbf{importo non positivo?} \textrightarrow{} \texttt{IllegalArgumentException}, non controllata: è una precondizione violata, il client la evita controllando l'importo (come fa \texttt{leggiImporto})
    \item \textbf{saldo insufficiente?} \textrightarrow{} un'eccezione nostra, \textbf{controllata} (\texttt{extends Exception}): dipende dallo stato del conto, il client deve prevederla e il compilatore glielo ricorda
    \item \textbf{che cosa porta con sé?} \textrightarrow{} quanto manca: una variabile di istanza \texttt{final} con il suo \texttt{get}; il messaggio, costruito nel costruttore, con \texttt{super(...)}
    \item \textbf{le intestazioni?} \textrightarrow{} \texttt{versa} e \texttt{preleva} sono \texttt{void}: l'esito non è più un \texttt{boolean}, chi chiama sa che l'operazione è riuscita se non arriva un'eccezione. Solo \texttt{preleva} e \texttt{trasferisci} dichiarano \texttt{throws}
    \item \textbf{in che ordine i controlli di \texttt{preleva}?} \textrightarrow{} prima l'importo, poi il saldo; il saldo si modifica solo dopo entrambi
    \item \textbf{\texttt{trasferisci}?} \textrightarrow{} chiama \texttt{preleva}, poi \texttt{altro.versa}: se \texttt{preleva} solleva l'eccezione, \texttt{trasferisci} termina lì e il versamento non avviene. Non la cattura: la rimanda al client, che sa che cosa fare
\end{enumerate}

\column{0.38\textwidth}
\MermaidAdjustBoxOpts{max width=\linewidth,max height=0.44\textheight,center}
\begin{mermaid}
classDiagram
  class ContoBancario {
    -intestatario : String
    -saldo : double
    +ContoBancario(String)
    +getIntestatario() String
    +getSaldo() double
    +versa(double)
    +preleva(double)
    +trasferisci(ContoBancario, double)
    +toString() String
  }
\end{mermaid}
\MermaidAdjustBoxOpts{max width=\linewidth,max height=0.2\textheight,center}
\begin{mermaid}
classDiagram
  class SaldoInsufficienteException {
    -mancante : double
    +SaldoInsufficienteException(double)
    +getMancante() double
  }
\end{mermaid}
\MermaidAdjustBoxOpts{max width=0.95\linewidth,max height=0.62\textheight,center}
\end{columns}
\end{frame}

% ========================================================================
\begin{frame}[fragile]{Una possibile soluzione (1/2) -- l'eccezione e l'inizio del conto}
\begin{columns}[T]
\column{0.49\textwidth}
\begin{lstlisting}[basicstyle=\ttfamily\scriptsize]
/** Sollevata se il saldo non basta. */
public class SaldoInsufficienteException
        extends Exception {

    private final double mancante;

    public SaldoInsufficienteException(
            double mancante) {
        super("saldo insufficiente, mancano "
                + mancante + " euro");
        this.mancante = mancante;
    }

    public double getMancante() {
        return mancante;
    }
}
\end{lstlisting}
\footnotesize
\begin{itemize}\itemsep=2pt
    \item \texttt{extends Exception}: controllata. Il costruttore riceve quanto manca, costruisce il messaggio per \texttt{getMessage()} e lo passa a \texttt{super(...)}, di norma la prima istruzione
    \item \texttt{mancante} è \texttt{final}: un'eccezione descrive un fatto avvenuto, non cambia
\end{itemize}

\column{0.51\textwidth}
\begin{lstlisting}[basicstyle=\ttfamily\scriptsize]
/** Un conto: errori con le eccezioni. */
public class ContoBancario {

    private final String intestatario;
    private double saldo;

    public ContoBancario(String intestatario) {
        this.intestatario = intestatario;
        saldo = 0;
    }

    public String getIntestatario() {
        return intestatario;
    }

    public double getSaldo() {
        return saldo;
    }
\end{lstlisting}
\end{columns}
\end{frame}

% ========================================================================
\begin{frame}[fragile]{Una possibile soluzione (2/2) -- \texttt{ContoBancario.java}}
\begin{columns}[T]
\column{0.5\textwidth}
\begin{lstlisting}[basicstyle=\ttfamily\scriptsize]
/** Versa importo, che deve essere > 0. */
public void versa(double importo) {
    if (importo <= 0) {
        throw new IllegalArgumentException(
                "non positivo: " + importo);
    }
    saldo += importo;
}

/** Preleva importo, se il saldo basta. */
public void preleva(double importo)
        throws SaldoInsufficienteException {
    if (importo <= 0) {
        throw new IllegalArgumentException(
                "non positivo: " + importo);
    }
    if (importo > saldo) {
        throw new SaldoInsufficienteException(
                importo - saldo);
    }
    saldo -= importo;
}
\end{lstlisting}

\column{0.5\textwidth}
\begin{lstlisting}[basicstyle=\ttfamily\scriptsize]
/** Sposta importo su altro. */
public void trasferisci(ContoBancario altro,
                        double importo)
        throws SaldoInsufficienteException {
    preleva(importo);
    altro.versa(importo);
}

@Override
public String toString() {
    return intestatario + ": saldo " + saldo;
}
\end{lstlisting}
\footnotesize
\begin{itemize}\itemsep=2pt
    \item in \texttt{preleva} i due \texttt{throw} lasciano il saldo com'era: si modifica solo all'ultima riga
    \item \texttt{versa} non dichiara nulla: \texttt{IllegalArgumentException} è non controllata
    \item \texttt{trasferisci} non cattura: se \texttt{preleva} la solleva, l'eccezione attraversa \texttt{trasferisci} e arriva al client; \texttt{altro.versa} non viene eseguito
\end{itemize}
\end{columns}
\end{frame}

% ========================================================================
\section{Wrap-up}

\begin{frame}{Riepilogo: cosa abbiamo trattato}
\footnotesize
\begin{center}
\renewcommand{\arraystretch}{1.1}
\scriptsize
\begin{tabular}{|>{\raggedright\arraybackslash}p{6.4cm}|>{\raggedright\arraybackslash}p{7.4cm}|}
\hline
\textbf{Concetto (unità)} & \textbf{Dove compare oggi} \\
\hline
Eccezioni a tempo di esecuzione, stack trace (U19) & \texttt{Divisione} con \texttt{7 0}; \texttt{Catena}; \texttt{SenzaBase} \\
\hline
\texttt{try} e \texttt{catch}, più \texttt{catch}, multi-catch, \texttt{finally} (U19) & \texttt{Divisione} v2, \texttt{DivisioneFinally}; \texttt{ProvaConto} \\
\hline
\texttt{throw}; stack unwinding; rilancio (U19) & \texttt{Media}; \texttt{Catena}, \texttt{CatenaRilancio}; \texttt{trasferisci} \\
\hline
Gerarchia; controllate e non; \texttt{throws} (U19) & \texttt{ProvaClone}; \texttt{preleva}, \texttt{trasferisci} \\
\hline
Le nostre eccezioni; quando non usare le eccezioni (U19) & \texttt{ParcheggioPienoException}, \texttt{SaldoInsufficienteException}; \texttt{Quoziente} \\
\hline
Input robusto con \texttt{Scanner} (U19) & \texttt{leggiImporto} in \texttt{ProvaConto} \\
\hline
Ricorsione: caso base, caso ricorsivo, convergenza (U21) & \texttt{fattoriale}; \texttt{sommaCifre}, \texttt{stampaInversa}, \texttt{fibonacci} \\
\hline
Lo stack delle chiamate; terminazione e profondità (U21) & il diagramma di \texttt{fattoriale(3)}; \texttt{StackOverflowError} \\
\hline
Ricorsione e iterazione; quando usare la ricorsione (U21) & \texttt{ContoAllaRovescia}; le 21\,891 chiamate di \texttt{fibonacci(20)} \\
\hline
\end{tabular}
\end{center}

\vspace{0.1cm}
\textbf{Prossima lezione (L11)}: la ricerca binaria, iterativa e ricorsiva, e quanto costa; la programmazione generica -- classi e metodi generici, vincoli con \texttt{extends}, interfacce generiche, il carattere jolly \texttt{?}.
\end{frame}

% ========================================================================
\begin{frame}{Homework 10 (da svolgere entro la prossima lezione)}
\footnotesize\setlength{\parskip}{0pt}
Tre esercizi, ciascuno con il suo programma di prova, nella cartella \texttt{homework/}. Nei primi due, niente cicli: solo ricorsione.
\begin{enumerate}\itemsep=3pt
    \item \textbf{\texttt{Hanoi.java}} (con \texttt{ProvaHanoi.java}): la \textbf{Torre di Hanoi}. Tre paletti \texttt{A}, \texttt{B}, \texttt{C}; su \texttt{A} \texttt{n} dischi dal più grande (in basso) al più piccolo; vanno portati su \texttt{C} spostando un disco alla volta, mai uno più grande sopra uno più piccolo. Per spostare \texttt{n} dischi da \texttt{A} a \texttt{C}: si spostano \texttt{n-1} dischi su \texttt{B}, il disco \texttt{n} su \texttt{C}, gli \texttt{n-1} da \texttt{B} su \texttt{C}. \texttt{Hanoi} è un modulo con il metodo \texttt{static} \texttt{sposta(int n, char da, char a, char appoggio)}, che stampa ogni mossa (\texttt{disco 1: A -> C}), e un contatore \texttt{private static} delle mosse con \texttt{getMosse()} e \texttt{azzera()}. \texttt{ProvaHanoi} legge \texttt{n} da \texttt{args[0]}: se l'argomento manca stampa come si usa (con un \texttt{if}); se non è un intero, \texttt{Integer.parseInt} solleva \texttt{NumberFormatException} (non controllata): catturatela. Se $n < 1$, un messaggio; altrimenti alla fine stampa le mosse e le attese, $2^n - 1$.
    \item \textbf{\texttt{Stringhe.java}} (con \texttt{ProvaStringhe.java}): un modulo con tre metodi \texttt{static} ricorsivi, che usano solo \texttt{length}, \texttt{charAt} e \texttt{substring}: \texttt{inverti(String s)} (\texttt{"ricorsione"} \textrightarrow{} \texttt{"enoisrocir"}), \texttt{isPalindroma(String s)} (primo e ultimo carattere uguali e la parte in mezzo palindroma) e \texttt{conta(char c, String s)} (quante volte \texttt{c} compare in \texttt{s}).
    \item \textbf{\texttt{Distributore.java}} (con \texttt{CreditoInsufficienteException.java}, \texttt{EsauritoException.java} e \texttt{ProvaDistributore.java}): un distributore di un solo prodotto: costruttore \texttt{Distributore(String prodotto, double prezzo, int pezzi)}, credito iniziale 0. \texttt{inserisci(double monete)} solleva \texttt{IllegalArgumentException} se le monete non sono positive; \texttt{eroga()} consegna un pezzo e scala il prezzo dal credito, oppure solleva \texttt{EsauritoException} (pezzi finiti, controllato per primo) o \texttt{CreditoInsufficienteException} (con \texttt{getMancante()}), entrambe controllate; \texttt{restituisci()} azzera il credito e lo restituisce; \texttt{toString}. \texttt{ProvaDistributore} ha un metodo \texttt{static} \texttt{provaErogazione(Distributore d)} che chiama \texttt{eroga} e gestisce le due eccezioni con due \texttt{catch} diversi; il \texttt{main} prova un credito insufficiente, due erogazioni riuscite, il prodotto esaurito, il resto e un inserimento negativo.
\end{enumerate}
\end{frame}

% ========================================================================
\begin{frame}
\begin{center}
{\Huge Domande?}\\[2em]
{\normalsize \href{mailto:mattia.fonisto@unina.it}{\texttt{mattia.fonisto@unina.it}}}
\end{center}
\end{frame}

\end{document}
